% GENERATED FILE. Edit canonical lessons, then run npm run generate:books.
% canonical-source-hash: fnv1a64:41c40da74bea0e53
% canonical-lessons: TE-C139-tadi, TE-C139-podi, TE-C139-peda, TE-C139-musali, TE-C139-prasiddha, TE-R139-first-pass-animals-the-body-and-more-things, TE-R139-second-pass-describing-words

\chapter{Wet, Dry, Poor, Old, Famous}
\label{ch:te-wet-dry-poor-old-famous}
\hlchaptermodality{139}

\begin{chapteropening}
\textbf{By the end of this chapter:} \emph{I can say wet, dry, poor, old and famous.}

Complete the last lesson of chapter 139: i can say wet, dry, poor, old and famous.
\end{chapteropening}

\section[taḍi]{\te{తడి} (taḍi) — wet}

\label{lesson:TE-C139-tadi}



\begin{quote}
Before the new one: say the Telugu for wide, then the Telugu for narrow.
\end{quote}

\subsection*{You'll want to know: \te{తడి}}
\textbf{\te{తడి}} — \emph{taḍi} — \textquotedblleft{}wet\textquotedblright{}.

\begin{grammarlens}[title={What you've built}]
One more word for this chapter.
\end{grammarlens}

\subsection*{Your turn}
\begin{itemize}
\raggedright
  \item \emph{Say it:} \emph{taḍi}
  \item \emph{Say it:} \emph{taḍi}, once more
  \item \emph{Say it:} say \emph{taḍi}
  \item \emph{Recall:} say the Telugu for wide, then the Telugu for narrow, then say \emph{taḍi} again
\end{itemize}

\begin{tcolorbox}[breakable,colback=teal!4,colframe=teal!35!black,title={Before you move on}]
What does \te{తడి} mean? (\textquotedblleft{}Wet\textquotedblright{}.) Say it once more.
\end{tcolorbox}
\section[poḍi]{\te{పొడి} (poḍi) — dry}

\label{lesson:TE-C139-podi}



\begin{quote}
Before the new one: say the Telugu for narrow, then the Telugu for wet.
\end{quote}

\subsection*{You'll want to know: \te{పొడి}}
\textbf{\te{పొడి}} — \emph{poḍi} — \textquotedblleft{}dry\textquotedblright{}.

\begin{grammarlens}[title={What you've built}]
One more word for this chapter.
\end{grammarlens}

\subsection*{Your turn}
\begin{itemize}
\raggedright
  \item \emph{Say it:} \emph{poḍi}
  \item \emph{Say it:} \emph{poḍi}, once more
  \item \emph{Say it:} say \emph{poḍi}
  \item \emph{Recall:} say the Telugu for narrow, then the Telugu for wet, then say \emph{poḍi} again
\end{itemize}

\begin{tcolorbox}[breakable,colback=teal!4,colframe=teal!35!black,title={Before you move on}]
What does \te{పొడి} mean? (\textquotedblleft{}Dry\textquotedblright{}.) Say it once more.
\end{tcolorbox}
\section[pēda]{\te{పేద} (pēda) — poor}

\label{lesson:TE-C139-peda}



\begin{quote}
Before the new one: say the Telugu for wet, then the Telugu for dry.
\end{quote}

\subsection*{You'll want to know: \te{పేద}}
\textbf{\te{పేద}} — \emph{pēda} — \textquotedblleft{}poor\textquotedblright{}.

\begin{grammarlens}[title={What you've built}]
One more word for this chapter.
\end{grammarlens}

\subsection*{Your turn}
\begin{itemize}
\raggedright
  \item \emph{Say it:} \emph{pēda}
  \item \emph{Say it:} \emph{pēda}, once more
  \item \emph{Say it:} say \emph{pēda}
  \item \emph{Recall:} say the Telugu for wet, then the Telugu for dry, then say \emph{pēda} again
\end{itemize}

\begin{tcolorbox}[breakable,colback=teal!4,colframe=teal!35!black,title={Before you move on}]
What does \te{పేద} mean? (\textquotedblleft{}Poor\textquotedblright{}.) Say it once more.
\end{tcolorbox}
\section[musali]{\te{ముసలి} (musali) — old (of people)}

\label{lesson:TE-C139-musali}



\begin{quote}
Before the new one: say the Telugu for dry, then the Telugu for poor.

Say \textbf{\te{తాత}} and \textbf{\te{అమ్మమ్మ}}, grandfather and grandmother.
\end{quote}

\subsection*{You'll want to know: \te{ముసలి}}
\textbf{\te{ముసలి}} — \emph{musali} — \textquotedblleft{}old (of people)\textquotedblright{}.

\begin{grammarlens}[title={What you've built}]
One more word for this chapter.
\end{grammarlens}

\subsection*{Your turn}
\begin{itemize}
\raggedright
  \item \emph{Say it:} \emph{musali}
  \item \emph{Say it:} \emph{musali}, once more
  \item \emph{Say it:} say \emph{musali}
  \item \emph{Recall:} say the Telugu for dry, then the Telugu for poor, then say \emph{musali} again
\end{itemize}

\begin{tcolorbox}[breakable,colback=teal!4,colframe=teal!35!black,title={Before you move on}]
What does \te{ముసలి} mean? (\textquotedblleft{}Old (of people)\textquotedblright{}.) Say it once more.
\end{tcolorbox}
\section[prasiddha]{\te{ప్రసిద్ధ} (prasiddha) — famous}

\label{lesson:TE-C139-prasiddha}



\begin{quote}
Before the new one: say the Telugu for poor, then the Telugu for old.
\end{quote}

\subsection*{You'll want to know: \te{ప్రసిద్ధ}}
\textbf{\te{ప్రసిద్ధ}} — \emph{prasiddha} — \textquotedblleft{}famous\textquotedblright{}.

\begin{grammarlens}[title={What you've built}]
One more word for this chapter.
\end{grammarlens}

\subsection*{Your turn}
\begin{itemize}
\raggedright
  \item \emph{Say it:} \emph{prasiddha}
  \item \emph{Say it:} \emph{prasiddha}, once more
  \item \emph{Say it:} say \emph{prasiddha}
  \item \emph{Recall:} say the Telugu for poor, then the Telugu for old, then say \emph{prasiddha} again
\end{itemize}

\begin{tcolorbox}[breakable,colback=teal!4,colframe=teal!35!black,title={Before you move on}]
What does \te{ప్రసిద్ధ} mean? (\textquotedblleft{}Famous\textquotedblright{}.) Say it once more.
\end{tcolorbox}
\section[(dialogue)]{Animals, the body and more things — a first pass}

\label{lesson:TE-R139-first-pass-animals-the-body-and-more-things}



\begin{quote}
Say the words for old and famous.
\end{quote}

\subsection*{Your turn}
Say these aloud:

\begin{itemize}
\raggedright
  \item a television, a computer, a phone, a ticket, a paisa
  \item a rupee, a bottle, a long skirt, a pocket, a belt
  \item socks, a coat, a dhoti, a blouse, a lotus
  \item clothes, a moustache, tears, a boat, a train
  \item a car, a bicycle, an aeroplane, grain, wheat
\end{itemize}

\begin{tcolorbox}[breakable,colback=teal!4,colframe=teal!35!black,title={Before you move on}]
A television? (\textbf{\te{టీవీ}}, \emph{ṭīvī}.) A computer? (\textbf{\te{కంప్యూటర్}}, \emph{kampyūṭar}.) A phone? (\textbf{\te{ఫోను}}, \emph{phōnu}.) A ticket? (\textbf{\te{టిక్కెట్టు}}, \emph{ṭikkeṭṭu}.) A paisa? (\textbf{\te{పైసా}}, \emph{paisā}.) A rupee? (\textbf{\te{రూపాయి}}, \emph{rūpāyi}.) A bottle? (\textbf{\te{బాటిలు}}, \emph{bāṭilu}.) A long skirt? (\textbf{\te{లంగా}}, \emph{laṅgā}.) A pocket? (\textbf{\te{జేబు}}, \emph{jēbu}.) A belt? (\textbf{\te{బెల్టు}}, \emph{beltu}.) Socks? (\textbf{\te{సాక్సు}}, \emph{sāksu}.) A coat? (\textbf{\te{కోటు}}, \emph{kōṭu}.) A dhoti? (\textbf{\te{ధోతి}}, \emph{dhōti}.) A blouse? (\textbf{\te{జాకెట్టు}}, \emph{jākeṭṭu}.) A lotus? (\textbf{\te{తామరపువ్వు}}, \emph{tāmarapuvvu}.) Clothes? (\textbf{\te{గుడ్డలు}}, \emph{guḍḍalu}.) A moustache? (\textbf{\te{మీసాలు}}, \emph{mīsālu}.) Tears? (\textbf{\te{కంటినీరు}}, \emph{kaṇṭinīru}.) A boat? (\textbf{\te{పడవ}}, \emph{paḍava}.) A train? (\textbf{\te{రైలు}}, \emph{railu}.) A car? (\textbf{\te{కారు}}, \emph{kāru}.) A bicycle? (\textbf{\te{సైకిలు}}, \emph{sāikilu}.) An aeroplane? (\textbf{\te{విమానం}}, \emph{vimānaṁ}.) Grain? (\textbf{\te{ధాన్యం}}, \emph{dhānyaṁ}.) Wheat? (\textbf{\te{గోధుమలు}}, \emph{gōdhumalu}.)
\end{tcolorbox}
\section[(dialogue)]{Describing words — a second pass}

\label{lesson:TE-R139-second-pass-describing-words}



\begin{quote}
Say the words for black pepper and cumin.
\end{quote}

\subsection*{Your turn}
Say these aloud:

\begin{itemize}
\raggedright
  \item black pepper, cumin, tamarind, a lemon, grapes
  \item a guava, sugar, tall, short, heavy
  \item light, soft, hard, clean, dirty
  \item fast; speed, beautiful, deep, wide, narrow
  \item wet, dry, poor, old, famous
\end{itemize}

\begin{tcolorbox}[breakable,colback=teal!4,colframe=teal!35!black,title={Before you move on}]
Black pepper? (\textbf{\te{మిరియాలు}}, \emph{miriyālu}.) Cumin? (\textbf{\te{జీలకర్ర}}, \emph{jīlakarra}.) Tamarind? (\textbf{\te{చింతపండు}}, \emph{cintapaṇḍu}.) A lemon? (\textbf{\te{నిమ్మకాయ}}, \emph{nimmakāya}.) Grapes? (\textbf{\te{ద్రాక్ష}}, \emph{drākṣa}.) A guava? (\textbf{\te{జామపండు}}, \emph{jāmapaṇḍu}.) Sugar? (\textbf{\te{పంచదార}}, \emph{pañcadāra}.) Tall? (\textbf{\te{పొడవు}}, \emph{poḍavu}.) Short? (\textbf{\te{పొట్టి}}, \emph{poṭṭi}.) Heavy? (\textbf{\te{బరువు}}, \emph{baruvu}.) Light? (\textbf{\te{తేలిక}}, \emph{tēlika}.) Soft? (\textbf{\te{మెత్తని}}, \emph{mettani}.) Hard? (\textbf{\te{గట్టి}}, \emph{gaṭṭi}.) Clean? (\textbf{\te{శుభ్రం}}, \emph{śubhraṁ}.) Dirty? (\textbf{\te{మురికి}}, \emph{murikī}.) Fast; speed? (\textbf{\te{వేగం}}, \emph{vēgaṁ}.) Beautiful? (\textbf{\te{అందమైన}}, \emph{andamaina}.) Deep? (\textbf{\te{లోతు}}, \emph{lōtu}.) Wide? (\textbf{\te{వెడల్పు}}, \emph{veḍalpu}.) Narrow? (\textbf{\te{ఇరుకు}}, \emph{iruku}.) Wet? (\textbf{\te{తడి}}, \emph{taḍi}.) Dry? (\textbf{\te{పొడి}}, \emph{poḍi}.) Poor? (\textbf{\te{పేద}}, \emph{pēda}.) Old? (\textbf{\te{ముసలి}}, \emph{musali}.) Famous? (\textbf{\te{ప్రసిద్ధ}}, \emph{prasiddha}.)
\end{tcolorbox}
